% paper.tex lines 42-90 (abstract and introduction up to the main theorem)
\begin{abstract}
We prove Barnette's conjecture: every finite simple cubic bipartite planar
\(3\)-vertex-connected graph has a Hamiltonian cycle.
\end{abstract}

\section{Introduction}

A graph is \emph{cubic} if every vertex has degree three, and
\emph{bipartite} if its vertices can be divided into two classes with
no edge inside either class. It is \emph{\(3\)-vertex-connected} if it
has at least four vertices and deleting any set of at most two vertices
leaves it connected. A \emph{Hamiltonian cycle} visits every vertex
exactly once. We consider finite simple undirected graphs throughout
the main statement.

Barnette's conjecture, recorded by Gr\"unbaum in the 1969 Waterloo
proceedings \cite{Barnette1969}, asserts that every cubic bipartite polyhedral
graph is Hamiltonian, where polyhedral means planar and
\(3\)-vertex-connected. It belongs to the line of questions initiated
by Tait's conjecture for cubic polyhedral graphs, disproved by
Tutte \cite{Tutte1946}; see \citet{APSW2016,GSW2023} for the history
and equivalent formulations.

Several partial results explain the role of face structure.
Goodey proved the case in which all faces are quadrilaterals or
hexagons \cite{Goodey1975}. Feder and Subi allowed arbitrarily
large faces in one class of a proper three-coloring of the faces,
requiring only quadrilaterals and hexagons in the other two classes
\cite{FederSubi2006}. Their argument uses proper quasi spanning trees of faces
after contracting a distinguished face class. Kardo\v{s} proved
the related Barnette--Goodey conjecture for cubic polyhedral graphs
with faces of size at most six, without assuming bipartiteness
\cite{Kardos2020}. More recently, Schnieders proved the
faces-at-most-eight case of Barnette's conjecture
\cite[Corollary~1.13]{Schnieders2025}. We remove the face-size
restriction within the full \(3\)-connected bipartite class.

\begin{theorem}\label{thm:main}
Every finite simple undirected graph that is cubic, bipartite,
planar, and \(3\)-vertex-connected has a Hamiltonian cycle.
\end{theorem}

The difficulty is to obtain one spanning cycle. A spanning subgraph
of degree two can have several cycle components, and local choices
alone need not join them. We instead construct an acyclic edge set
in a dual triangulation whose complement becomes one connected
spanning cycle. The construction also retains a prescribed local
pattern needed when the triangulation is cut into smaller pieces.

